\chapter{Proposed solution}
\label{cap:cap3}

\section{System requirements and High-Level Design}

\subsection{Problem Reiteration and Solution Strategy}

In the current digital health landscape, nutritional tracking applications frequently fail 
to sustain long-term user adherence due to structural rigidity. Traditional platforms 
rely on manual, repetitive data logging where users must explicitly search, weigh, and 
log every individual food item consumed. Furthermore, generic meal planning engines often 
deliver static, pre-packaged templates that fail to account for unique metabolic shifts, 
individual lifestyle activity factors, or lean muscle mass preservation requirements.

The strategy proposed in this thesis addresses these core limitations by shifting the 
paradigm from manual tracking to automated macro-nutrient synthesis. This is achieved 
through a decentralized microservice architecture that decouples user identity tracking 
from mathematical plan generation. The solution introduces an autonomous optimization 
engine, the OptimizerService, which dynamically calculates nutritional targets and 
automatically compiles meal configurations tailored to the user's specific metabolic 
profile. By automating the data-entry layer and providing instant, scientifically backend 
plan variations, the platform eliminates the traditional points of friction in dietary 
logging while maintaining high structural flexibility.


\subsection{Original Ideas and Innovative Solutions}
\begin{itemize}
    \item \textbf{Embedded Read-Only Database for Meal Optimization:} The system separates data storage concerns across two distinct database engines based on their access patterns. Transactional data such as user profiles, authentication state, and daily plans are persisted in a cloud-hosted Azure SQL instance. The \texttt{OptimizerService}, which handles meal generation, operates exclusively on static pre-seeded nutritional data that never changes at runtime. For this reason, an embedded SQLite database is deployed directly alongside the service container, eliminating network round-trips to a cloud database for read-only lookups during meal plan generation. Additionally, the \texttt{DailyPlanService} maintains local snapshots of this nutritional data, removing any runtime dependency on the \texttt{OptimizerService}'s internal storage.

    \item \textbf{Dynamic Lean Mass Preservation Target Calibration:} Unlike typical platforms that apply rigid calorie deficits across all profiles, this solution implements a dynamically scaled allocation algorithm within the \texttt{DailyPlanService}. The system evaluates the user's active configuration against specific constraints using the following relationship:
    
    \begin{equation}
        TDEE = BMR \times \text{ActivityFactor}
    \end{equation}
    
    When a restrictive deficit is selected, the system automatically recalibrates macro-nutrient distributions, scaling protein allocations up to 1.4g/kg or 2.0g/kg based on calculated lean mass preservation metrics while strictly enforcing a minimum fat threshold of 0.7g/kg or 0.8g/kg to protect metabolic health. This separation ensures nutritional logic remains within the .NET domain layer, keeping the \texttt{OptimizerService} focused solely on meal selection.

    \item \textbf{Stateless Authentication with Refresh Token Rotation:} The system implements a dual-token authentication strategy to balance security and user experience across a cross-domain microservices architecture. Short-lived JWT access tokens (15 minutes) are stored in JavaScript memory and transmitted via the Authorization Bearer header, eliminating CSRF attack vectors entirely. Long-lived refresh tokens (30 days) are persisted server-side in the \texttt{AuthServiceDb} and transmitted exclusively via HttpOnly cookies, making them inaccessible to JavaScript and therefore resistant to XSS attacks. Upon each token refresh, the existing refresh token is deleted and replaced with a newly generated one — a practice known as refresh token rotation — ensuring that a stolen refresh token can only be used once before being invalidated.
\end{itemize}

\section{User Requirements}

The application's scope is driven by the practical needs of the end user, divided into functional operations and non-functional quality expectations.

\subsection{Functional Requirements}
\begin{itemize}
    \item \textbf{Authentication:} Users must be able to sign up, log in securely, and 
    remain authenticated across page refreshes.
    \item \textbf{Onboarding Form:} First-time users require a guided form to enter basic 
    metrics (weight, height, age, sex, and activity level).
    \item \textbf{Metabolic Calculation:} The system must automatically compute and display 
    the user's BMR and TDEE based on their onboarding input.
    \item \textbf{Diet Goal Customization:} Users need to select a specific goal 
    (e.g., Lean Cut, Lean Bulk, or Lose Fat) which automatically scales their calorie and 
    macro targets.
    \item \textbf{Meal Plan Selection:} Users must be able to choose between automatically 
    generating a meal plan via the AI optimizer or manually adding foods.
    \item \textbf{Dashboard CRUD:} The main interface must allow users to view their daily 
    macro progress, log consumed food items, update metrics, and delete entries.
\end{itemize}

\subsection{Non-Functional Requirements}
\begin{itemize}
    \item \textbf{Performance:} Page transitions, database saves, and macro calculations 
    must happen within a sub-second timeframe to ensure a smooth user experience.
    \item \textbf{Data Privacy:} User profiles, target metrics, and logged food records must 
    be securely isolated so users can only access their own data.
\end{itemize}

\section{Architectural Modeling}

\subsection{System Architecture}

The system is designed using a decentralized microservice architecture to ensure 
scalability, fault isolation, and strict separation of business concerns. Each service 
operates independently with its own data store, exposing functionality exclusively through 
well-defined HTTP APIs. The structural organization of the system's components, their 
interfaces, and data boundaries are illustrated in the component diagram below.

\begin{figure}[htbp]
    \centering
    \includegraphics[width=0.9\textwidth]{continut/capitol3/figuri/ComponentDiagram.pdf}
    \caption{System Component Architecture Diagram}
    \label{fig:component_architecture}
\end{figure}

As shown in Figure~\ref{fig:component_architecture}, the system architecture is composed of 
the following core layers and components:

\begin{itemize}
    \item \textbf{Presentation Layer (React SPA):} A lightweight, client-side single-page 
    application built with React. It handles user interactions, manages local routing states, 
    and communicates asynchronously with the backend APIs via standard RESTful HTTP protocols 
    using JSON payloads.
    \item \textbf{User Management Service (\texttt{UserManagementService}):} A specialized 
    backend service responsible for user administration. It processes user registration, 
    validates credentials during login, handles secure session management, and stores physical 
    user profile metrics.
    \item \textbf{Daily Plan Service (\texttt{DailyPlanService}):} The orchestrator for 
    the user's active nutritional tracking. It handles core CRUD operations related to food
    logging, manages daily macro-nutrient targets based on selected goals, and coordinates
    meal plan persistence. When automated meal generation is requested, it delegates to the
    \texttt{OptimizerService} via a typed \texttt{HttpClient}, forwarding the user's 
    targets and receiving back an optimized meal configuration.
    \item \textbf{Optimizer Service (\texttt{OptimizerService}):} A highly isolated 
    computational engine built exclusively for running the meal generation algorithms. It 
    ingests targeted macro goals and selects optimal food combinations from its embedded 
    storage without interfering with transactional user flows.
    \item \textbf{Data Tier Topology:} Storage concerns are explicitly segregated. Relational 
    transactional records and user authorization states are persisted in the 
    cloud-hosted \texttt{AuthServiceDb} (Azure SQL). In contrast, the static food and 
    ingredient catalogs used for macro synthesis are hosted within an embedded SQLite 
    database running locally inside the \texttt{OptimizerService} container to minimize 
    network latency.
\end{itemize}

\clearpage
\subsection{Database Design}

The system employs three distinct databases, each scoped to its respective service to 
ensure strict data isolation and separation of concerns.

\subsubsection{UserManagement Service Database}
The \textit{AuthServiceDb} manages user identity, authentication state, and physical profile 
data. The \textit{AspNetUsers} table is provided by the ASP.NET Core Identity framework and 
serves as the central user entity. The \textit{UserDetails} table maintains a one-to-one 
relationship with \textit{AspNetUsers}, storing user physiological data used for 
nutritional calculations. The \textit{RefreshTokens} table maintains a one-to-many 
relationship with \textit{AspNetUsers}, persisting refresh tokens issued after successful 
authentication to support token rotation.

\begin{figure}[htbp]
    \centering
    \includegraphics[width=0.9\textwidth]{continut/capitol3/figuri/AuthServiceERdiagram.pdf}
    \caption{UserManagement Service Database Schema}
    \label{fig:auth_db}
\end{figure}

\clearpage
\subsubsection{DailyPlan Service Database}
The \textit{DailyPlanServiceDb} stores user meal plans and their associated nutritional 
data. Each \textit{DailyPlan} record contains one or more \textit{PlannedMeals}, which 
in turn contain one or more \textit{PlannedMealIngredients}. This structure represents a 
snapshot of the generated meal plan, persisted independently from the OptimizerService to 
avoid runtime coupling between services.

\begin{figure}[htbp]
    \centering
    \includegraphics[width=0.9\textwidth]{continut/capitol3/figuri/DailyPlanServiceERdiagram.pdf}
    \caption{DailyPlan Service Database Schema}
    \label{fig:dailyplan_db}
\end{figure}

\clearpage
\subsubsection{Optimizer Service Database}
The OptimizerService uses an embedded SQLite database pre-seeded with static nutritional 
data. The \textit{basic\_foods} table stores individual food items with their macronutrient 
values. The \textit{meals} table defines meal templates with complexity and type 
classifications. The \textit{meal\_ingredients} table serves as a junction table 
establishing a many-to-many relationship between meals and foods. This database is
 read-only at runtime and is bundled directly within the service container.

\begin{figure}[htbp]
    \centering
    \includegraphics[width=0.9\textwidth]{continut/capitol3/figuri/OptimizerServiceERdiagram.pdf}
    \caption{Optimizer Service Database Schema}
    \label{fig:optimizer_db}
\end{figure}

\section{Implementation Description}
This section describes the primary classes and programming components 
developed to implement the platform's core business logic. The backend 
architecture follows a clean, layered structure where controllers expose 
the API endpoints, services handle business operations, and entities 
model the core data structures.

\subsection{Description of Implemented Elements and Classes}

Both backend services follow a uniform three-layer structure: a controller at the API
boundary, a service class containing the domain logic, and a \texttt{DbContext} injected
directly into the service for data access. This pattern is applied consistently across
all features, with each layer depending on its neighbour exclusively through an interface,
enabling independent testability and clear separation of concerns. The components described
below represent the most significant classes implemented within this structure.

\begin{itemize}
    \item \textbf{\texttt{UserController}:} Positioned within the API exposure
    layer of the \texttt{UserManagementService}. It exposes the complete HTTP
    interface for session and profile lifecycle management, covering both
    anonymous operations such as registration, login, and token refresh, and
    authenticated operations such as logout, identity retrieval, and physical
    metric persistence. Security boundaries are enforced at the controller level
    using the \texttt{[Authorize]} attribute on all routes that require a valid
    access token.
    
    \item \textbf{\texttt{UserDetailsService}:} Implements the core domain logic for user 
    processing. It coordinates payload validation, computes basic metabolic baselines 
    from raw physical data, and triggers persistent repository writes to the database 
    tier.
    
    \item \textbf{\texttt{TokenService}:} A dedicated utility component within 
    the authentication infrastructure that handles the cryptographic generation 
    of short-lived JSON Web Tokens (JWT). It populates the token payload with 
    essential user identity properties using standard security claims, specifically 
    mapping the unique user identifier (\texttt{ClaimTypes.NameIdentifier}) and the 
    validated email address (\texttt{ClaimTypes.Email}). The resulting JWT is 
    cryptographically signed using the \texttt{HmacSha256Signature} algorithm and 
    issued with a strict 15-minute expiration lifespan.

    \item \textbf{\texttt{RefreshTokenService}:} Implements the \texttt{IRefreshTokenService} 
    interface to handle the lifecycle of long-lived session tokens. It provides three 
    core operations: \texttt{GenerateAndStoreAsync} to cryptographically generate a 
    secure token string and persist it in the database for a specific user, 
    \texttt{ValidateAsync} to retrieve and check the validity of an incoming token
    against storage, and \texttt{DeleteAsync} to remove a token upon user logout, 
    enforcing the refresh token rotation strategy.

    \item \textbf{\texttt{IdentityService}:} Implements the \texttt{IIdentityService} 
    interface to manage core user account operations. It abstracts user interaction 
    with the security provider through three asynchronous tasks: \texttt{AuthenticateAsync} 
    to validate plaintext user credentials against stored password hashes, 
    \texttt{RegisterAsync} to handle the creation of new user records within the 
    database, and \texttt{FindByIdAsync} to look up an active user profile using their 
    unique identifier.

    \item \textbf{\texttt{DailyPlanController}:} Exposes the API endpoints for managing 
    tracking records and triggering menu synthesis. It handles the integration between 
    the user front-end and the optimization core through two key 
    operations: \texttt{POST /api/dailyplan/generate}, which redirects user physical 
    parameters to the internal optimization engine to build a macro-calibrated menu, 
    and \texttt{GET /api/dailyplan/ingredient/search}, which queries available food 
    data pools. Additionally, it implements core CRUD operations across both the meal 
    and ingredient levels, specifically allowing users to fetch complete plans, modify 
    specific meal titles, and add, update, or delete individual ingredients within a 
    planned meal slot.
    
    \item \textbf{\texttt{DailyPlanService}:} Implements the \texttt{IDailyPlanService} 
    interface to execute the business logic behind the active tracking dashboard. 
    It acts as the direct coordinator between the data repositories and the controller 
    layer, handling state modifications across both meal and ingredient scopes. It 
    processes date-specific plan lookups (\texttt{GetDailyPlanAsync}), manages full 
    layout initializations or structural replacements, and evaluates modification rules 
    whenever users add, update, or remove ingredients and meal slot parameters.
    
    \item \textbf{\texttt{MealOptimizerClient}:} Implements the
    \texttt{IMealOptimizerClient} interface to abstract the communication with
    the underlying optimization engine service. It provides a clean, async gateway
    to forward user metabolic configurations and macro-nutrient targets down to the
    solver (\texttt{OptimizeAsync}) and also forwards the keyword-matching
    logic (\texttt{GetIngredientOptionsAsync}) to populate available food selections
    during custom plan creation.

    \item \textbf{\texttt{FitTrackWithMLP.Shared}:} A dedicated shared class library
    referenced by both the \newline\texttt{UserManagementService} and the
    \texttt{DailyPlanService}. It centralizes cross-cutting concerns that would
    otherwise be duplicated across services. It defines the common Data Transfer
    Objects used for inter-layer communication, such as user profile payloads and
    plan request structures, as well as shared enumerations including goal type
    statuses and plan generation states. Beyond contracts, the library also houses
    a global exception handling middleware applied uniformly across both services,
    and an \texttt{AddFitTrackAuthentication} extension method on
    \texttt{IServiceCollection} that encapsulates the shared JWT validation
    configuration, ensuring both services enforce identical token verification
    rules from a single definition.
\end{itemize}

\subsection{End-to-End Request Flow}

\subsection{Algorithm Pseudocode}

\section{Technology Stack Selection and Justification}

\subsection{Hardware Platform Analysis}

\subsection{Advantages and Disadvantages}

\section{Application Walkthrough and Screenshots}